El hardware utilizado para realizar los experimentos corresponde a un procesador Intel Core i5-12450H, 16 GB de memoria RAM DDR4 a 3200 MT/s y sistema operativo Windows 11. Los algoritmos fueron ejecutados en un entorno Ubuntu mediante WSL2, el cual tenía asignado aproximadamente 7,6 GiB de memoria RAM. Para la compilación de los programas en C++ se utilizó g++ (GCC) versión 13.3.0, utilizando el estándar C++17.

Los datasets utilizados corresponden a los archivos generados mediante grid\_generator.py, siguiendo las instrucciones y configuraciones definidas para la tarea. Se utilizaron grillas de dimensiones N×M, considerando valores de N iguales a 4, 8, 16, 64, 256 y 1024, mientras que M podía tomar los valores N o 2N. Además, se consideraron cuatro tipos de grillas: libres, con murallas, con minas y mixtas, utilizando los dominios D10 y D1000, con tres muestras por configuración. En total se generaron 288 instancias utilizando la semilla 2026 para permitir que, cada vez que se ejecute el generador, se obtengan los mismos casos de prueba. Los tres algoritmos fueron evaluados sobre las instancias generadas, con la excepción de Fuerza Bruta, que se ejecutó únicamente cuando N+M era menor o igual a 32, siguiendo las restricciones establecidas en la tarea.
